\chapter{Physical and Theoretical Framework}

This chapter explains the physical model behind the operator hierarchy. The
later chapters are intentionally abstract, but the abstraction is not detached
from the grid. Every matrix in the thesis comes from a physical layer: phasor
voltages, line admittances, active-power balance, rotor or converter dynamics,
controller states, and the operating point around which the system is
linearized. The purpose of this chapter is to make that chain explicit before
the thesis uses \(L\), \(M\), \(D\), \(\Pi(s)\), and \(T(s)^{-1}\) as technical
objects.

\thesisbox{thesisblue}{Chapter message}{
The problem is not only to compute eigenvalues. The problem is to identify
which physical mechanism creates, moves, amplifies, or protects those
eigenvalues. That is why the thesis needs a hierarchy from network stiffness to
fixed damping, rational self-energy, Green functions, structured fragility, and
nonlinear memory.}

\section{How to Read This Chapter Step by Step}

The rest of the thesis uses compact notation, so this chapter starts from the
physical chain and builds the operator one piece at a time. The intended
reading order is:
\begin{enumerate}
\item Start with the electrical graph: buses are nodes, lines are edges, and
line admittances create the network equation \(I=Y_{\rm bus}V\).
\item Fix an operating point: the same topology can have different
small-signal stiffness under a different dispatch, voltage profile, or angle
stress.
\item Linearize active-power balance: the derivative of power with respect to
angles becomes the synchronizing stiffness matrix \(L\).
\item Add dynamic storage and damping: inertia gives \(M\ddot\theta\), damping
or damping-like control gives \(D\dot\theta\), and stiffness gives
\(L\theta\).
\item Introduce the spectral parameter \(s\): an oscillatory response
\(\theta(t)=ve^{st}\) converts the differential equation into an operator
\(s^2M+sD+L\).
\item Escalate only when needed: if hidden controller states are eliminated,
the missing frequency-dependent physics appears as \(\Pi(s)\), and the object
becomes \(T(s)=s^2M+sD_0+L+\Pi(s)\).
\end{enumerate}
The key point is that each symbol is a physical commitment. A graph eigenvalue
is not automatically a stability certificate; it is a certificate only when
the dynamic terms that can change the protected margin have been accounted for.

\section{The Grid as a Networked Dynamic System}

A power grid is a networked dynamical system, not a single aggregate machine.
The buses form the electrical graph, the lines and transformers define the
admittance matrix, the generators and converters inject dynamic behavior, and
the loads close the power balance. A useful small-signal model must preserve
all four roles: network geometry, operating-point physics, dynamic storage, and
control memory.

\begin{figure}[ht]
\centering
\resizebox{0.96\textwidth}{!}{%
\begin{tikzpicture}[
  bus/.style={circle,draw=thesisink!65,fill=white,line width=0.45pt,minimum size=6pt,inner sep=0pt},
  gen/.style={circle,draw=thesisblue!85!black,fill=thesisblue!10,line width=0.9pt,minimum size=8pt,inner sep=0pt},
  conv/.style={regular polygon,regular polygon sides=4,draw=thesisorange!85!black,fill=thesisorange!12,line width=0.9pt,minimum size=9pt,inner sep=0pt},
  load/.style={rectangle,draw=thesisgray!85!black,fill=softgray,line width=0.6pt,minimum width=8pt,minimum height=6pt,inner sep=0pt},
  line/.style={draw=thesisgray!76,line width=0.7pt},
  ctrl/.style={rectangle,rounded corners=2pt,draw=thesisred!75!black,fill=softred,align=center,font=\scriptsize,text width=2.25cm},
  lbl/.style={font=\scriptsize\color{thesisink}},
  flow/.style={-{Latex[length=2.0mm]},draw=thesisgray,line width=0.55pt}
]
\coordinate (b1) at (0,1.6);
\coordinate (b2) at (1.7,2.2);
\coordinate (b3) at (3.4,1.55);
\coordinate (b4) at (5.0,2.15);
\coordinate (b5) at (6.5,1.25);
\coordinate (b6) at (1.1,0.25);
\coordinate (b7) at (2.8,-0.35);
\coordinate (b8) at (4.8,0.05);
\coordinate (b9) at (6.2,-0.55);
\draw[line] (b1)--(b2)--(b3)--(b4)--(b5)--(b9)--(b8)--(b7)--(b6)--(b1);
\draw[line] (b2)--(b6) (b2)--(b7) (b3)--(b7) (b3)--(b8) (b4)--(b8) (b5)--(b8);
\foreach \p in {b1,b4,b6,b8} \node[bus] at (\p) {};
\node[gen] (g1) at (b2) {};
\node[conv] (c1) at (b3) {};
\node[gen] (g2) at (b5) {};
\node[conv] (c2) at (b7) {};
\node[load] (l1) at (b9) {};
\node[lbl,above=0.08cm of g1] {SG};
\node[lbl,above=0.08cm of c1] {IBR};
\node[lbl,above=0.08cm of g2] {SG};
\node[lbl,below=0.10cm of c2] {GFM/GFL};
\node[lbl,below=0.08cm of l1] {load};
\node[ctrl] (pll) at (3.4,3.20) {PLL/current loop\\voltage loop\\plant control};
\node[ctrl] (gov) at (0.95,3.20) {governor\\exciter\\stabilizer};
\draw[flow] (gov.south) -- (g1.north);
\draw[flow] (pll.south) -- (c1.north);
\node[lbl,align=center,text width=4.2cm] at (9.1,1.80) {\textbf{Electrical graph}\\
\(Y_{\rm bus}=G+jB\)\\
\(I=Y_{\rm bus}V\)};
\node[lbl,align=center,text width=4.2cm] at (9.1,0.60) {\textbf{Dynamic devices}\\
\(M\ddot\theta+D\dot\theta\)\\
or converter states \(z\)};
\node[lbl,align=center,text width=4.2cm] at (9.1,-0.65) {\textbf{Planning actions}\\
line, damping, inertia,\\
controller retuning};
\draw[flow] (6.95,1.25) -- (7.65,1.62);
\draw[flow] (6.95,0.18) -- (7.65,0.48);
\draw[flow] (6.95,-0.48) -- (7.65,-0.62);
\end{tikzpicture}}
\caption{A converter-rich grid contains a physical graph, dynamic devices, and
hidden controller states. The retained variables may be bus or machine angles,
frequencies, voltages, and selected ports; the eliminated variables reappear as
self-energy when they are condensed.}
\label{fig:framework_networked_grid}
\end{figure}

Figure~\ref{fig:framework_networked_grid} shows the modeling issue. The
electrical graph tells us which buses are connected, but not which dynamic
margin is protected. A synchronous generator contributes inertia, damping,
exciter, governor, and stabilizer dynamics. A grid-following or grid-forming
converter contributes synchronization, current-loop, voltage-loop, filter,
limiter, and plant-level dynamics. If those internal states are removed from a
reduced model, their effect must go somewhere. In this thesis, that place is
the rational self-energy \(\Pi(s)\).

\section{AC Phasors, Power Flow, and Network Stiffness}

The first physical object is the AC phasor voltage. At bus \(i\),
\begin{equation}
  V_i=|V_i|e^{j\theta_i},
  \label{eq:framework_phasor}
\end{equation}
where \(|V_i|\) is voltage magnitude and \(\theta_i\) is phase angle. The
network admittance matrix maps bus voltages to currents,
\begin{equation}
  I=Y_{\rm bus}V,\qquad Y_{\rm bus}=G+jB,
  \label{eq:framework_ybus}
\end{equation}
and the complex power injected at the bus is
\begin{equation}
  S_i=P_i+jQ_i=V_iI_i^* .
  \label{eq:framework_complex_power}
\end{equation}
The operating point is therefore not only a graph. It is a solution of the AC
power-flow equations at a particular dispatch, load, voltage profile, topology,
and controller setting.

The active-power stiffness matrix is the first derivative of active power with
respect to phase angles at the operating point. With voltage magnitudes fixed
inside the small-signal angle block,
\begin{equation}
  \Delta P\approx J_{P\theta}\Delta\theta,
  \qquad
  J_{P\theta}
  =
  \left.
  \frac{\partial P}{\partial \theta}
  \right|_{(V^\star,\theta^\star)} .
  \label{eq:framework_jacobian}
\end{equation}
Depending on the sign convention for exported electrical power, the
synchronizing matrix is \(L=J_{P\theta}\) or \(L=-J_{P\theta}\). The thesis
always states the convention before using a derivative. The important fact is
that \(L\) is an operating-point derivative, not a purely topological matrix.

\begin{figure}[ht]
\centering
\resizebox{0.92\textwidth}{!}{%
\begin{tikzpicture}[
  axis/.style={-{Latex[length=2.0mm]},draw=thesisink!70,line width=0.55pt},
  curve/.style={draw=thesisblue,line width=1.05pt},
  tangent/.style={draw=thesisorange,line width=0.85pt},
  guide/.style={draw=thesisgray!70,densely dashed,line width=0.45pt},
  lbl/.style={font=\small\color{thesisink}}
]
\draw[axis] (-3.55,0) -- (3.65,0) node[right,lbl] {\(\delta=\theta_i-\theta_j\)};
\draw[axis] (0,-1.75) -- (0,1.95) node[above,lbl] {\(P_{ij}\)};
\draw[curve] plot[smooth] coordinates {
  (-3.14,0) (-2.62,-0.50) (-2.09,-0.87) (-1.57,-1.00)
  (-1.05,-0.87) (-0.52,-0.50) (0,0) (0.52,0.50)
  (1.05,0.87) (1.57,1.00) (2.09,0.87) (2.62,0.50) (3.14,0)
};
\coordinate (op) at (0.78,0.70);
\draw[guide] (op) -- (0.78,0) node[below,lbl] {\(\delta^\star\)};
\draw[guide] (op) -- (0,0.70);
\draw[tangent] (-0.35,-0.26) -- (1.85,1.66);
\node[circle,fill=thesisred,inner sep=1.7pt] at (op) {};
\node[lbl,anchor=west] at (1.00,0.44) {operating point};
\node[lbl,align=center,text width=4.8cm] at (-2.05,1.32)
{\(\displaystyle P_{ij}\approx |V_i||V_j|B_{ij}\sin\delta\)};
\node[lbl,align=center,text width=5.3cm] at (2.05,-1.25)
{\(\displaystyle \frac{\partial P_{ij}}{\partial \delta}
=|V_i^\star||V_j^\star|B_{ij}\cos\delta^\star\)};
\end{tikzpicture}}
\caption{A line contributes synchronizing stiffness because active power is
sensitive to angle separation. The slope at the operating point becomes the
edge weight in the linearized stiffness matrix \(L\).}
\label{fig:framework_power_angle_plot}
\end{figure}

Figure~\ref{fig:framework_power_angle_plot} gives the physical reason \(L\)
acts like a stiffness matrix. For a lossless two-bus line,
\begin{equation}
  P_{ij}\approx |V_i||V_j|B_{ij}\sin(\theta_i-\theta_j).
  \label{eq:framework_line_power}
\end{equation}
Linearization gives a slope proportional to
\(|V_i^\star||V_j^\star|B_{ij}\cos(\theta_i^\star-\theta_j^\star)\). This
slope is large when the line is electrically strong and the operating point is
not close to the angular stability limit. It decreases as the operating point
approaches a stressed angle difference. Thus the same physical line can have
different small-signal stiffness under different operating conditions.

\section{From Physical Balance to the Swing Equation}

The second physical object is rotational or emulated energy storage. For a
synchronous machine, the rotor accelerates when mechanical input and
electrical output are imbalanced:
\begin{equation}
  M_i\dot\omega_i=P_{m,i}-P_{e,i}-D_i\omega_i,
  \qquad
  \omega_i=\dot\theta_i .
  \label{eq:framework_single_swing}
\end{equation}
After linearization around a balanced operating point,
\begin{equation}
  M_i\Delta\ddot\theta_i+D_i\Delta\dot\theta_i+\Delta P_{e,i}=p_i .
  \label{eq:framework_single_linear_swing}
\end{equation}
The network supplies \(\Delta P_e\approx L\Delta\theta\). Stacking retained
angle variables gives
\begin{equation}
  M\ddot\theta+D\dot\theta+L\theta=p .
  \label{eq:framework_network_swing}
\end{equation}
The three matrices are not interchangeable. \(M\) multiplies acceleration,
\(D\) multiplies velocity, and \(L\) multiplies angle displacement. A planning
action that changes one term should not be interpreted as if it changed the
others.

The uniform angle direction has no restoring force. Mathematically,
\begin{equation}
  L\ones=0 .
  \label{eq:framework_reference_invariance}
\end{equation}
Physically, adding the same constant angle to every bus does not change any
angle difference or power flow. All spectral statements therefore occur after
fixing an angle reference, using center-of-inertia coordinates, or projecting
onto the nonzero-angle subspace.

\section{Converter States and the DAE Layer}

Converter-rich grids add a hidden dynamic layer between the electrical graph
and the retained swing-like variables. A generic local model has differential
states \(x\), algebraic variables \(y\), and parameters or actions \(\rho\):
\begin{align}
  \dot x &= f(x,y,\rho), \label{eq:framework_dae_f}\\
  0 &= g(x,y,\rho). \label{eq:framework_dae_g}
\end{align}
The algebraic variables include network constraints such as nodal current
balance and voltage equations. The differential states include machine rotor
states, exciter states, governor states, PLL states, current controllers,
voltage controllers, filters, and plant controls.

At a regular equilibrium, \(g_y\) is nonsingular and the algebraic variables
can be eliminated locally. The first-order small-signal matrix is
\begin{equation}
  A_{\rm red}=f_x-f_y g_y^{-1}g_x .
  \label{eq:framework_reduced_A}
\end{equation}
This formula explains why the operating point matters. The Jacobians \(f_x\),
\(f_y\), \(g_x\), and \(g_y\) are all evaluated at
\((x^\star,y^\star,\rho^\star)\). If the dispatch, IBR penetration, line
status, voltage setpoints, or controller gains change, the linearized operator
changes.

\begin{figure}[ht]
\centering
\resizebox{\textwidth}{!}{%
\begin{tikzpicture}[
  layer/.style={rectangle,rounded corners=3pt,draw=#1!75!black,fill=#1!9,
    minimum height=1.0cm,text width=2.85cm,align=center,font=\small},
  arrow/.style={-{Latex[length=2.2mm]},draw=thesisgray,line width=0.7pt},
  note/.style={font=\scriptsize\color{thesisgray},align=center,text width=2.8cm}
]
\node[layer=thesisblue] (net) {Electrical network\\\(Y_{\rm bus},V,\theta\)};
\node[layer=thesisgreen, right=0.55cm of net] (opf) {Operating point\\\(P,Q,V^\star,\theta^\star\)};
\node[layer=thesisorange, right=0.55cm of opf] (jac) {Linearization\\\(L,J_{xy},A_{\rm red}\)};
\node[layer=thesisred, right=0.55cm of jac] (dyn) {Dynamic operator\\\(T(s)\)};
\node[layer=thesisteal, right=0.55cm of dyn] (dec) {Decision\\margin, action,\\certificate};
\draw[arrow] (net) -- (opf);
\draw[arrow] (opf) -- (jac);
\draw[arrow] (jac) -- (dyn);
\draw[arrow] (dyn) -- (dec);
\node[note, below=0.25cm of net] {topology and admittance};
\node[note, below=0.25cm of opf] {dispatch and constraints};
\node[note, below=0.25cm of jac] {small-signal physics};
\node[note, below=0.25cm of dyn] {QEP or Schur NEP};
\node[note, below=0.25cm of dec] {hosting or retuning};
\end{tikzpicture}}
\caption{The thesis follows the physical modeling chain from network data to
an operator-valued decision. A planning result is not reproducible unless the
operating point, retained variables, action family, margin, and certification
gate are all declared.}
\label{fig:framework_modeling_pipeline}
\end{figure}

\section{Why the Spectral Parameter Appears}

Small-signal stability studies the free response after an infinitesimal
disturbance. For the fixed-\(D\) second-order model, set \(p=0\) and try
\(\theta(t)=ve^{st}\). Since
\[
  \dot\theta=sve^{st},
  \qquad
  \ddot\theta=s^2ve^{st},
\]
substitution into \eqref{eq:framework_network_swing} gives
\begin{equation}
  (s^2M+sD+L)v=0.
  \label{eq:framework_qep}
\end{equation}
This quadratic eigenvalue problem is the first spectral object of the thesis.
The term \(s^2M\) is inertia, the term \(sD\) is velocity damping, and the
term \(L\) is synchronizing stiffness.

When controller states are eliminated rather than retained, the operator is no
longer a fixed-coefficient QEP. Partition the linearized dynamics into retained
states \(x_r\) and condensed states \(x_c\). Eliminating \(x_c\) gives a Schur
complement with the resolvent \((sI-A_{cc})^{-1}\). In second-order retained
coordinates this appears as
\begin{equation}
  T(s)=s^2M+sD_0+L+\Pi(s),
  \qquad
  \Pi(s)=\Pi_K(s)+s\Pi_D(s).
  \label{eq:framework_nep}
\end{equation}
The rational term \(\Pi(s)\) is the spectral memory of the eliminated
controller block. It contains poles and residues. A static damping matrix
cannot exactly represent those poles.

\section{What a Pole Means Physically}

A root \(s=\sigma+j\omega\) of the retained operator produces a modal response
of the form
\[
  e^{st}=e^{\sigma t}e^{j\omega t}.
\]
The real part \(\sigma\) controls growth or decay, and the imaginary part
\(\omega\) controls oscillation frequency. The damping ratio is
\begin{equation}
  \zeta=-\frac{\sigma}{\sqrt{\sigma^2+\omega^2}} .
  \label{eq:framework_damping_ratio}
\end{equation}
The same pole can be interpreted in several physical ways: a slow
electromechanical oscillation, a local plant mode, a controller interaction
mode, or a mixed network-control mode. This is why the thesis does not name a
mode only by frequency. It asks which operator term and which port direction
make the mode visible.

\begin{figure}[ht]
\centering
\resizebox{0.86\textwidth}{!}{%
\begin{tikzpicture}[
  axis/.style={-{Latex[length=2.0mm]},draw=thesisink!75,line width=0.55pt},
  stable/.style={fill=softgreen,draw=none,opacity=0.55},
  protect/.style={draw=thesisorange,line width=0.85pt,densely dashed},
  pole/.style={circle,fill=thesisblue,inner sep=1.7pt},
  badpole/.style={circle,fill=thesisred,inner sep=1.8pt},
  arrow/.style={-{Latex[length=2.1mm]},draw=thesisred,line width=0.8pt},
  lbl/.style={font=\small\color{thesisink}}
]
\fill[stable] (-4.0,-2.0) rectangle (0,2.25);
\draw[axis] (-4.05,0) -- (1.45,0) node[right,lbl] {\(\Real(s)\)};
\draw[axis] (0,-2.05) -- (0,2.35) node[above,lbl] {\(\Imag(s)\)};
\draw[protect] (-1.0,-2.0) -- (-1.0,2.15);
\node[lbl,anchor=east] at (-1.05,1.92) {protected boundary};
\node[pole] (p1) at (-2.15,1.18) {};
\node[pole] (p2) at (-2.15,-1.18) {};
\node[pole] (p3) at (-0.58,0.78) {};
\node[badpole] (p4) at (0.18,0.72) {};
\draw[arrow] (p3) -- (p4);
\node[lbl,align=center,text width=3.2cm] at (-2.75,-1.70) {stable and protected};
\node[lbl,align=center,text width=3.4cm] at (0.52,1.38) {action or uncertainty can move a pole};
\end{tikzpicture}}
\caption{The spectral plane makes the margin explicit. A nominally stable pole
can still be fragile if a physically admissible perturbation moves it across a
protected damping, abscissa, resonance, or passivity boundary.}
\label{fig:framework_spectral_plane}
\end{figure}

Figure~\ref{fig:framework_spectral_plane} explains why a nominal eigenvalue
plot is not enough. A system may be stable at the base point but fragile under
small line uncertainty, controller retuning, model-identification error, or
operating-point movement. The thesis therefore uses structured fragility
rather than only nominal roots.

\section{Why Scalar Grid Strength Is Not Enough}

Classical strength metrics are useful because they compress a complicated
network into a small number. Short-circuit ratio describes local electrical
stiffness. Algebraic connectivity describes a global graph bottleneck. A
damping ratio describes one pole. Effective resistance describes a static
graph compliance. None of these scalars is wrong. The problem is that each one
protects only one mechanism.

\begin{table}[ht]
\centering
\caption{Why a scalar metric is not a complete dynamic-strength theory.}
\label{tab:framework_scalar_limits}
\begin{tabular}{L{0.20\textwidth}L{0.28\textwidth}L{0.34\textwidth}}
\toprule
\rowcolor{softblue}
\thead{Scalar} & \thead{What it sees} & \thead{What it can miss}\\
\midrule
SCR & local fault-level or Thevenin stiffness & damping, controller poles,
non-normal amplification, multi-port interactions\\
\(\lambda_2(L)\) & global graph connectivity & heterogeneous damping,
inertia placement, operating-point residues\\
\(\zeta_k\) & decay of one pole & finite-action fragility, nearby controller
poles, hidden input-output gain\\
\(a_e^TL^\dagger a_e\) & static effective resistance & frequency-dependent
compliance, resonant Green-function peaks\\
\bottomrule
\end{tabular}
\end{table}

The dynamic replacement is not another scalar. It is the retained operator
\begin{equation}
  T(s)=s^2M+sD_0+L+\Pi(s)
  \label{eq:framework_core_operator}
\end{equation}
and its Green function
\begin{equation}
  G_T(s)=T(s)^{-1}.
  \label{eq:framework_green}
\end{equation}
The scalar quantities used for planning are projections of this operator onto a
declared port, action family, frequency band, or protected boundary.

\begin{figure}[ht]
\centering
\resizebox{0.86\textwidth}{!}{%
\begin{tikzpicture}[
  axis/.style={-{Latex[length=2.0mm]},draw=thesisink!75,line width=0.55pt},
  curveA/.style={draw=thesisblue,line width=1.0pt},
  curveB/.style={draw=thesisred,line width=1.0pt},
  curveC/.style={draw=thesisorange,line width=1.0pt,densely dashed},
  lbl/.style={font=\small\color{thesisink}},
  leg/.style={font=\scriptsize\color{thesisink},anchor=west}
]
\draw[axis] (0,0) -- (7.55,0) node[right,lbl] {line stiffness \(b_e\)};
\draw[axis] (0,0) -- (0,3.55) node[above,lbl] {normalized margin};
\draw[curveA] plot[smooth] coordinates {(0.45,0.55) (1.15,1.02) (1.95,1.55) (2.85,2.04) (3.75,2.38) (4.70,2.66) (5.55,2.84)};
\draw[curveB] plot[smooth] coordinates {(0.45,2.80) (1.15,2.60) (1.95,2.30) (2.85,1.88) (3.75,1.38) (4.70,0.92) (5.55,0.62)};
\draw[curveC] plot[smooth] coordinates {(0.45,2.45) (1.15,2.42) (1.95,2.26) (2.85,1.90) (3.75,1.20) (4.70,0.78) (5.55,1.02)};
\draw[curveA] (5.95,2.70) -- (6.45,2.70);
\node[leg] at (6.55,2.70) {graph stiffness};
\draw[curveB] (5.95,2.33) -- (6.45,2.33);
\node[leg] at (6.55,2.33) {damping headroom};
\draw[curveC] (5.95,1.96) -- (6.45,1.96);
\node[leg] at (6.55,1.96) {controller-aware margin};
\end{tikzpicture}}
\caption{A stiffness-improving line action need not improve a dynamic margin.
The graph metric, fixed-\(D\) damping screen, and controller-aware margin can
move in different directions.}
\label{fig:framework_scalar_failure_plot}
\end{figure}

Figure~\ref{fig:framework_scalar_failure_plot} is the central planning
motivation. If a line upgrade increases graph stiffness but reduces damping
headroom, a scalar graph score gives the wrong engineering message. If the
static screen accepts the line but the controller-aware model rejects it, the
missing mechanism is not topology. It is controller memory.

\section{A Three-Bus Example Where Classical Screens Fail}

This section gives the smallest possible example that contains the dissertation
problem. The example is not meant to represent a full test system. It is a
teaching model that shows why a line action that is excellent for a classical
graph metric can be bad for a dynamic margin.

\subsection{Step 1: Build the Three-Bus Graph}

Consider the three-bus chain in Figure~\ref{fig:framework_three_bus}. Bus 1 is
a synchronous source with appreciable damping, bus 2 is a passive transfer
bus, and bus 3 is a converter-dominated bus with weak damping. The base system
has two lines, \(1\)--\(2\) and \(2\)--\(3\), both with unit synchronizing
weight. The candidate planning action is a new tie line \(1\)--\(3\) with
incremental stiffness \(p\ge 0\).

\begin{figure}[ht]
\centering
\resizebox{0.90\textwidth}{!}{%
\begin{tikzpicture}[
  bus/.style={circle,draw=thesisink!70,fill=white,line width=0.7pt,
    minimum size=10mm,align=center,font=\small},
  gen/.style={circle,draw=thesisblue!85!black,fill=thesisblue!12,line width=1.0pt,
    minimum size=11mm,align=center,font=\small},
  conv/.style={regular polygon,regular polygon sides=4,draw=thesisorange!85!black,
    fill=thesisorange!13,line width=1.0pt,minimum size=11mm,align=center,font=\small},
  line/.style={draw=thesisgray!78,line width=1.2pt},
  cand/.style={draw=thesisred!82,line width=1.0pt,densely dashed},
  lbl/.style={font=\small\color{thesisink},align=center},
  note/.style={rectangle,rounded corners=2pt,draw=thesisgray!45,fill=softgray,
    align=center,font=\scriptsize,text width=3.3cm}
]
\node[gen] (b1) at (0,0) {1};
\node[bus] (b2) at (3.8,0) {2};
\node[conv,rotate=45] (b3) at (7.6,0) {\rotatebox{-45}{3}};
\draw[line] (b1) -- node[above,lbl] {\(b_{12}=1\)} (b2);
\draw[line] (b2) -- node[above,lbl] {\(b_{23}=1\)} (b3);
\draw[cand] (b1) to[bend right=60] node[pos=0.50,below=0.70cm,lbl] {candidate tie \(p\)} (b3);
\node[lbl,below=0.36cm of b1] {SG\\\(d_1=0.60\)};
\node[lbl,below=0.36cm of b2] {transfer bus\\\(d_2=0.60\)};
\node[lbl,below=0.36cm of b3] {IBR bus\\\(d_3=0.02\)};
\node[note] at (1.45,1.60) {classical graph screen:\\increase connectivity};
\node[note] at (6.05,1.60) {dynamic screen:\\protect damping and resonance};
\end{tikzpicture}}
\caption{Three-bus teaching example. The candidate tie \(1\)--\(3\) improves
the static graph, but it also changes the frequency of the endpoint mode that
uses the weakly damped converter bus.}
\label{fig:framework_three_bus}
\end{figure}

\subsection{Step 2: Derive the Stiffness Matrix}

With only the chain lines present, the Laplacian stiffness is
\begin{equation}
  L_0=
  \begin{bmatrix}
  1 & -1 & 0\\
  -1 & 2 & -1\\
  0 & -1 & 1
  \end{bmatrix}.
  \label{eq:framework_three_bus_L0}
\end{equation}
The nonzero graph modes are
\begin{equation}
  \nu_1=1,\quad
  \phi_1=\frac{1}{\sqrt 2}
  \begin{bmatrix}1\\0\\-1\end{bmatrix},
  \qquad
  \nu_2=3,\quad
  \phi_2=\frac{1}{\sqrt 6}
  \begin{bmatrix}1\\-2\\1\end{bmatrix}.
  \label{eq:framework_three_bus_modes}
\end{equation}
The candidate tie has incidence vector
\begin{equation}
  a_{13}=e_1-e_3=
  \begin{bmatrix}1\\0\\-1\end{bmatrix},
  \qquad
  L(p)=L_0+p\,a_{13}a_{13}^T .
  \label{eq:framework_three_bus_action}
\end{equation}
Because \(a_{13}\) is aligned with \(\phi_1\), the tie line changes the first
nonzero stiffness exactly:
\begin{equation}
  \nu_1(p)=1+2p,\qquad \nu_2(p)=3 .
  \label{eq:framework_three_bus_eigs}
\end{equation}
For \(0\le p\le 1\), the algebraic connectivity therefore increases from
\(1\) to \(3\). The static effective resistance between buses 1 and 3 also
decreases:
\begin{equation}
  R_{13}^{\rm stat}(p)=\frac{2}{1+2p}.
  \label{eq:framework_three_bus_resistance}
\end{equation}
A classical graph-only decision would say that increasing \(p\) is
unambiguously good: the network becomes more connected and the endpoint
resistance becomes smaller.

\subsection{Step 3: Add the First Dynamic Screen}

Now add unit inertia and heterogeneous damping,
\begin{equation}
  M=I,\qquad
  D=\operatorname{diag}(0.60,0.60,0.02).
  \label{eq:framework_three_bus_MD}
\end{equation}
The endpoint mode \(\phi_1\) moves bus 1 against bus 3 while bus 2 stays nearly
still. Its frozen graph-mode damping estimate is
\begin{equation}
  \delta_1=\phi_1^TD\phi_1
  =\frac{d_1+d_3}{2}
  =0.31 .
  \label{eq:framework_three_bus_modal_damping}
\end{equation}
The same line that increases stiffness also raises the modal frequency. A
first dynamic damping-ratio screen gives
\begin{equation}
  \zeta_1(p)\approx
  \frac{\delta_1}{2\sqrt{\nu_1(p)}}
  =
  \frac{0.31}{2\sqrt{1+2p}} .
  \label{eq:framework_three_bus_zeta}
\end{equation}
The sign of the conclusion changes. The graph metric improves with \(p\), but
the damping ratio falls with \(p\). If the planning rule requires
\(\zeta\ge 0.10\), the base case passes and the strong tie \(p=1\) fails.

\begin{figure}[ht]
\centering
\resizebox{0.82\textwidth}{!}{%
\begin{tikzpicture}[
  axis/.style={-{Latex[length=2.0mm]},draw=thesisink!75,line width=0.55pt},
  grid/.style={draw=thesisgray!24,line width=0.35pt},
  stiff/.style={draw=thesisblue,line width=1.1pt},
  damp/.style={draw=thesisred,line width=1.1pt},
  ctrl/.style={draw=thesisorange,line width=1.0pt,densely dashed},
  thresh/.style={draw=thesisink!55,line width=0.6pt,densely dotted},
  lbl/.style={font=\small\color{thesisink}},
  smalllbl/.style={font=\scriptsize\color{thesisink}},
  leg/.style={font=\scriptsize\color{thesisink},anchor=west}
]
\foreach \x in {0.70,1.95,3.20,4.45,5.70} \draw[grid] (\x,0.70) -- (\x,3.35);
\foreach \y in {0.95,1.50,2.05,2.60,3.15} \draw[grid] (0.55,\y) -- (5.95,\y);
\draw[axis] (0.55,0.70) -- (6.25,0.70) node[right,lbl] {tie strength \(p\)};
\draw[axis] (0.55,0.70) -- (0.55,3.55) node[above,lbl] {normalized score};
\node[smalllbl] at (0.70,0.43) {0};
\node[smalllbl] at (3.20,0.43) {0.5};
\node[smalllbl] at (5.70,0.43) {1};
\draw[stiff] plot[smooth] coordinates {(0.70,1.05) (1.95,1.52) (3.20,2.03) (4.45,2.55) (5.70,3.05)};
\draw[damp] plot[smooth] coordinates {(0.70,2.92) (1.95,2.42) (3.20,1.85) (4.45,1.42) (5.70,1.18)};
\draw[ctrl] plot[smooth] coordinates {(0.70,2.65) (1.95,2.06) (3.20,1.33) (4.45,0.88) (5.70,0.80)};
\draw[thresh] (0.55,1.33) -- (5.95,1.33);
\node[smalllbl,anchor=east] at (0.46,1.33) {requirement};
\draw[stiff] (0.90,0.08) -- (1.32,0.08);
\node[leg] at (1.42,0.08) {\(\lambda_2(L(p))\)};
\draw[damp] (2.80,0.08) -- (3.22,0.08);
\node[leg] at (3.32,0.08) {\(\zeta_1(p)\)};
\draw[ctrl] (4.45,0.08) -- (4.87,0.08);
\node[leg] at (4.97,0.08) {controller distance};
\end{tikzpicture}}
\caption{In the three-bus example, the classical graph score and the dynamic
margin move in opposite directions. This is the simplest visible failure of a
connectivity-only planning rule.}
\label{fig:framework_three_bus_plot}
\end{figure}

\subsection{Step 4: Show the Classical Failure}

Table~\ref{tab:framework_three_bus_failure} summarizes the contradiction. The
line action is attractive under two classical static metrics:
\(\lambda_2(L(p))\) rises and \(R_{13}^{\rm stat}(p)\) falls. The same action
is unattractive under the dynamic damping screen because it raises the
frequency of a mode that uses the weakly damped converter bus.

\begin{table}[ht]
\centering
\small
\caption{Three-bus example: static improvement but dynamic degradation.}
\label{tab:framework_three_bus_failure}
\begin{tabular}{c c c c c L{0.23\textwidth}}
\toprule
\rowcolor{softblue}
\thead{\(p\)} & \thead{\(\lambda_2\)} & \thead{\(R_{13}^{\rm stat}\)}
& \thead{\(\omega_1\approx\sqrt{1+2p}\)}
& \thead{\(\zeta_1\)} & \thead{Decision if \(\zeta_{\rm req}=0.10\)}\\
\midrule
0 & 1.00 & 2.00 & 1.00 & 0.155 & passes the damping screen\\
0.5 & 2.00 & 1.00 & 1.41 & 0.110 & barely passes\\
1.0 & 3.00 & 0.67 & 1.73 & 0.089 & fails even though graph strength improved\\
\bottomrule
\end{tabular}
\end{table}

This is the basic reason the thesis cannot stop at graph strength. A graph
method answers, ``Did the network become stiffer?'' The stability question is
different: ``Did the protected dynamic margin improve?'' In this example the
answer to the first question is yes, while the answer to the second is no.

\subsection{Step 5: Add the Hidden Controller Mechanism}

The previous step used only \(M\), \(D\), and \(L(p)\). Converter-rich systems
can fail one level higher. Suppose the eliminated converter/controller block at
bus 3 has a lightly damped internal pole near
\begin{equation}
  \mu_c=-0.05+j1.70 .
  \label{eq:framework_three_bus_controller_pole}
\end{equation}
A static graph screen cannot see this pole. A fixed-\(D\) screen cannot
represent it exactly either, because the reduced effect is frequency
dependent. After Schur condensation, it appears inside \(\Pi(s)\), and the
mode must be tested through
\begin{equation}
  T(s;p)=s^2I+sD+L(p)+\Pi(s).
  \label{eq:framework_three_bus_nep}
\end{equation}
A simple resonance-distance gate is
\begin{equation}
  d_c(p)=
  \frac{|j\sqrt{1+2p}-\mu_c|}
       {|j\sqrt{1+2p}|+|\mu_c|}.
  \label{eq:framework_three_bus_controller_distance}
\end{equation}
For \(p=0\), \(d_c\approx0.260\). For \(p=0.5\), \(d_c\approx0.093\). For
\(p=1\), \(d_c\approx0.017\). Thus the same tie that the classical graph
method ranks best also moves the network mode toward the hidden controller
pole. The correct escalation is not cosmetic; it changes the engineering
decision. The controller-aware analysis asks whether
\(\Pi(s)\), residues, port direction, and dynamic effective resistance create a
fragile interaction near the protected frequency band.

\subsection{What the Example Teaches}

The three-bus example separates three statements that are often conflated:
\begin{enumerate}
\item \textbf{Topology statement:} adding the tie \(1\)--\(3\) increases
\(\lambda_2\) and lowers static resistance.
\item \textbf{Fixed-dynamics statement:} the same tie can lower damping ratio
because it raises frequency without adding proportional damping.
\item \textbf{Controller-aware statement:} if a hidden controller pole sits in
the new frequency band, the reduced operator must include \(\Pi(s)\), and the
planning question becomes a port-resolvent question.
\end{enumerate}
This is why the dissertation develops a hierarchy instead of a single graph
index. Classical methods are valuable first screens, but they fail when the
protected margin depends on heterogeneous damping, hidden controller memory, or
finite line actions. The rest of the thesis formalizes exactly when each level
is enough and when the analysis must escalate.

\section{Dynamic Effective Resistance as the Bridge}

The classical effective resistance of a line \(e=(i,j)\) is
\begin{equation}
  R_e^{\rm stat}=a_e^TL^\dagger a_e,
  \qquad a_e=e_i-e_j .
  \label{eq:framework_static_resistance}
\end{equation}
This quantity measures static graph compliance between the endpoints of the
line. The thesis generalizes it by replacing the static pseudoinverse with the
dynamic Green function:
\begin{equation}
  \mathcal R_e(s)=a_e^TT(s)^{-1}a_e .
  \label{eq:framework_dynamic_resistance}
\end{equation}
At \(s=0\) and \(\Pi(0)=0\), this reduces to
\(a_e^TL^\dagger a_e\). Away from zero, it is a frequency-dependent,
controller-aware port compliance.

This definition matters because a finite line action has an exact scalar pole
condition:
\begin{equation}
  \det\!\left(T(s)+\Delta b_ea_ea_e^T\right)
  =
  \det T(s)\left(1+\Delta b_e\mathcal R_e(s)\right).
  \label{eq:framework_line_determinant}
\end{equation}
Thus the new or shifted poles created by the line action satisfy
\begin{equation}
  1+\Delta b_e\mathcal R_e(s)=0.
  \label{eq:framework_line_condition}
\end{equation}
The usual line sensitivity is the local residue limit of this finite-action
law. That is the reason dynamic effective resistance is the bridge between
spectral graph theory and controller-aware planning.

\section{Operator Hierarchy and Escalation Logic}

The thesis uses a hierarchy because no single model level is always necessary
or always sufficient. The cheapest valid model should be used, but it must pass
the gate attached to the claim.

\begin{figure}[ht]
\centering
\resizebox{\textwidth}{!}{%
\begin{tikzpicture}[
  level/.style={rectangle,rounded corners=3pt,draw=#1!75!black,fill=#1!9,
    minimum height=1.12cm,text width=2.7cm,align=center,font=\small},
  arrow/.style={-{Latex[length=2.2mm]},draw=thesisgray,line width=0.7pt},
  gate/.style={diamond,aspect=2.1,draw=thesisink!55,fill=softyellow,
    align=center,font=\scriptsize,text width=2.15cm,inner sep=1pt}
]
\node[level=thesisblue] (l0) {\(D=\gamma M\)\\uniform decay};
\node[gate, right=0.40cm of l0] (g0) {commuting?};
\node[level=thesisgreen, right=0.40cm of g0] (l1) {fixed \(D\)\\QEP};
\node[gate, right=0.40cm of l1] (g1) {modal coupling small?};
\node[level=thesisorange, right=0.40cm of g1] (l2) {Schur NEP\\\(\Pi(s)\)};
\node[gate, below=0.70cm of l2] (g2) {finite action or uncertainty?};
\node[level=thesisred, left=0.55cm of g2] (l3) {Green function\\\(T(s)^{-1}\)};
\node[level=thesisteal, right=0.55cm of g2] (l4) {structured fragility\\\(r_{\rm frag}\)};
\node[level=thesisgray, below=0.70cm of g2] (l5) {nonlinear memory\\\(K(x,\tau)\)};
\draw[arrow] (l0) -- (g0);
\draw[arrow] (g0) -- (l1);
\draw[arrow] (l1) -- (g1);
\draw[arrow] (g1) -- (l2);
\draw[arrow] (l2) -- (g2);
\draw[arrow] (g2) -- (l3);
\draw[arrow] (g2) -- (l4);
\draw[arrow] (g2) -- (l5);
\end{tikzpicture}}
\caption{Escalation logic of the thesis. The analysis becomes richer only
when the protected claim requires the richer operator: modal coupling,
controller memory, finite actions, structured uncertainty, or nonlinear
amplitude effects.}
\label{fig:framework_escalation}
\end{figure}

Figure~\ref{fig:framework_escalation} is the practical answer to why the full
analysis is necessary. If damping commutes with stiffness, graph modes can
certify more. If damping does not commute, the fixed-\(D\) QEP is needed. If
controller states are condensed, \(\Pi(s)\) is needed. If the question concerns
a finite line, damping, inertia, or controller action, \(T(s)^{-1}\) is needed.
If the question concerns minimum physical perturbation to failure, structured
fragility is needed. If the question concerns finite-amplitude behavior,
nonlinear memory is needed.

\section{Decision Contract for the Rest of the Thesis}

Every claim in the rest of the dissertation should be readable through the
same four-part contract:
\begin{enumerate}
\item \textbf{Operator:} Which model is being used: \(L\), QEP, Schur NEP,
Green function, structured fragility, or nonlinear memory?
\item \textbf{Port or action:} Which line, bus, converter, controller channel,
input-output pair, or uncertainty family is being tested?
\item \textbf{Margin:} Which quantity is protected: damping ratio, spectral
abscissa, protected gain, passivity, resonance distance, or hosting feasibility?
\item \textbf{Certificate:} Which theorem, bound, contour test, finite-action
law, or full-model validation turns the calculation into evidence?
\end{enumerate}
Without these four items, a weak-bus list, line-ranking plot, or hosting number
is not reproducible. With these four items, the thesis can distinguish a graph
screen from a dynamic certificate.

\FloatBarrier
